%! Equations and Graphs in Both Directions — Unit 2, Lesson 2.3 — Homework
\documentclass[10pt]{article}
\usepackage{atda-article}
\usepackage{atda-boxes}
\newcommand{\TallMath}[1]{$\displaystyle #1\rule[-1.4em]{0pt}{3.2em}$}
\setlength{\workrowsep}{8pt}

\begin{document}

\pageheader{Unit 2, Lesson 2.3}{Homework}
% NAMESTRIP: no \namedateperiod here. The student writes their name once, on the
% cover this component is stapled behind. See references/conventions.md.

\vspace{0.15in}

\boxguard[24]
\begin{notesbox}{Practice --- Both Directions}
\small
Half of these go equation $\to$ graph and half go graph $\to$ equation. Every equation you write
gets a substitution check.

\begin{enumerate}[leftmargin=1.6em, itemsep=6pt, topsep=4pt]

\item \textbf{Anchor points.} Give the anchor point of each image. For the exponential, give the
equation of its flat line (asymptote) as well.
\quad (a) $y=(x-6)^2+2$ \quad (b) $y=2^x+4$ \quad (c) $y=(x+1)^2-5$

\writelines{2}

\item \textbf{Equation $\to$ graph.} For $y=(x-3)^2-2$, complete the table, then plot the five points
on the grid and join them. The parent is drawn.
\smallskip

\renewcommand{\arraystretch}{1.25}
\begin{tabular}{@{}|c|c|c|c|c|c|@{}}
\hline
$x$ & $1$ & $2$ & $3$ & $4$ & $5$ \\ \hline
$y=(x-3)^2-2$ & \blank{1.0cm} & \blank{1.0cm} & \blank{1.0cm} & \blank{1.0cm} & \blank{1.0cm} \\
\hline
\end{tabular}

\begin{center}
\begin{tikzpicture}
\begin{axis}[
    width=9.2cm, height=4.6cm,
    xlabel={$x$}, ylabel={$y$},
    xmin=-3.4, xmax=6.4, ymin=-3.4, ymax=5.4,
    xtick={-3,-2,-1,0,1,2,3,4,5,6}, ytick={-3,-2,-1,0,1,2,3,4,5},
    grid=both, grid style={linegray!45}, axis lines=left,
    tick label style={font=\tiny}, label style={font=\scriptsize},
    legend entries={{$y=x^2$ (the parent)}},
    legend style={font=\scriptsize, draw=none, fill=none,
      at={(0.5,1.02)}, anchor=south, legend columns=2}]
  \addplot[royal, very thick, domain=-2.2:2.2, samples=80]{x^2};
\end{axis}
\end{tikzpicture}
\end{center}

\item \textbf{Graph $\to$ equation, in words.} Write the equation of each function described.
\begin{enumerate}[label=(\alph*), leftmargin=1.6em, itemsep=3pt, topsep=3pt]
  \item A parabola with exactly the parent's shape, lowest point at $(-4,3)$.
  \item An exponential whose asymptote is $y=5$ and which passes through $(0,6)$.
  \item A line through $(0,-1)$ and $(4,7)$.
\end{enumerate}

\writelines{3}

\item \textbf{Finding $a$.} A parabola has vertex $(1,2)$ and passes through $(2,5)$. Find $a$, write
the equation in transformation form, then check it at $x=3$.
\begin{work}
  5 &= a(2-1)^2+2 \\
  5 &= a+2 \\
  3 &= a
\end{work}

\writelines{2}

\item \textbf{Keystrokes.} Match each thing typed into a graphing calculator with what its screen
actually shows. Write the letter next to each description.
\smallskip

\renewcommand{\arraystretch}{1.35}
\begin{tabularx}{\linewidth}{@{}p{3.0cm} X@{}}
\rowcolor{mist}
\textbf{Typed} & \textbf{What the screen shows} \\
\hline
(a) \texttt{2\^{}X-3} & \blank{1.0cm} the parent $y=x^2$, unchanged \\
(b) \texttt{2\^{}(X-3)} & \blank{1.0cm} $y=2^x$ slid down $3$ \\
(c) \texttt{-X\^{}2} & \blank{1.0cm} $y=2^x$ slid right $3$ \\
(d) \texttt{(-X)\^{}2} & \blank{1.0cm} $y=x^2$ flipped over the $x$-axis \\
\end{tabularx}

\end{enumerate}
\end{notesbox}

\vspace{0.10in}

\boxguard[30]
\begin{scenariobox}[In Context --- The Yearbook's Profit Curve]{royal}
\small
The yearbook staff has to pick a price. The graph below came from three years of sales records:
$P(x)$ is the yearbook's profit, in dollars, when the price is $x$ dollars. The staff wants the
profit at a price of \$35, and \$35 is not marked.

\begin{center}
\begin{tikzpicture}
\begin{axis}[
    width=10.6cm, height=4.6cm,
    xlabel={$x$ (price of a yearbook, dollars)}, ylabel={profit (dollars)},
    xmin=0, xmax=62, ymin=0, ymax=2000,
    xtick={0,10,20,30,40,50,60},
    ytick={0,200,400,600,800,1000,1200,1400,1600,1800,2000},
    grid=both, grid style={linegray!45}, minor x tick num=1,
    minor grid style={linegray!30}, axis lines=left,
    tick label style={font=\tiny}, label style={font=\scriptsize},
    legend entries={{$P$ (profit from the sales records)}},
    legend style={font=\scriptsize, draw=none, fill=none,
      at={(0.5,1.02)}, anchor=south, legend columns=1}]
  \addplot[royal, very thick, domain=0:60, samples=100]{-2*(x-30)^2+1800};
\end{axis}
\end{tikzpicture}
\end{center}

\begin{enumerate}[leftmargin=1.6em, itemsep=6pt, topsep=2pt, start=6]

\item Read the vertex and both points where the curve meets the horizontal axis. Starting from the
parent $f(x)=x^2$, name the two translations and the reflection that put the curve there.

\writelines{2}

\item The shape is not the parent's, so find $a$. Use the right-hand point where the curve meets the
horizontal axis, $(60,0)$.
\begin{work}
  0 &= a(60-30)^2+1800 \\
  0 &= 900a+1800 \\
  -1800 &= 900a \\
  -2 &= a
\end{work}

\item Write $P(x)$ in transformation form, then answer the staff's question: what is the profit at a
price of \$35? Show the substitution.
\begin{work}
  P(x) &= -2(x-30)^2+1800 \\
  P(35) &= -2(35-30)^2+1800 \\
        &= -2(25)+1800 \\
        &= 1750
\end{work}

\item Say why the graph alone could not have given you item 8's answer. Then: what would you type
into the calculator, what window would you set, and what must the \textsc{table} show at
$\texttt{X}=35$? Finally, which of the domain and the range did the $+1800$ move?

\writelines{4}

\end{enumerate}
\end{scenariobox}

\vspace{0.10in}

\boxguard[24]
\begin{scenariobox}[A First Look Ahead]{goldacc}
\small
\begin{enumerate}[leftmargin=1.6em, itemsep=6pt, topsep=2pt, start=10]

\item Use the yearbook graph above --- no new equation needed. Answer in plain English where you do
not have the vocabulary yet; Lesson 2.4 will give you the words.
\begin{enumerate}[label=(\alph*), leftmargin=1.6em, itemsep=5pt, topsep=4pt]
  \item At which two prices is the profit exactly \$0? (These inputs have a name coming next
        lesson.)

\writelines{1}

  \item What is the highest profit the staff can reach, and at what price does it happen? Give
        \emph{both} numbers --- a value and a location.

\writelines{1}

  \item Over which prices is the profit \emph{going up} as the price goes up, and over which is it
        \emph{going down}? Describe the two stretches of prices.

\writelines{2}
\end{enumerate}

\end{enumerate}
\end{scenariobox}

\vspace{0.10in}

\boxguard[26]
\begin{extensionbox}
\small
\textbf{A stretch, a flip, and two slides at once --- plus a window.} Consider
$y=-2(x-4)^2+8$.

\begin{enumerate}[label=(\alph*), leftmargin=1.6em, itemsep=5pt, topsep=4pt]
  \item Name the parent and all three things done to it. Give the vertex, and say whether the
        parabola is narrower or wider than the parent.

\writelines{2}

  \item Find the two inputs where $y=0$ by solving. Then give four window values that would show
        both of them and the vertex.
\begin{work}
  0 &= -2(x-4)^2+8 \\
  -8 &= -2(x-4)^2 \\
  4 &= (x-4)^2 \\
  x-4 &= \pm 2 \\
  x &= 2 \text{ or } x = 6
\end{work}

\writelines{1}

% Unguarded, the extension box fit whole on p3 and pushed the remindbox ALONE onto
% p4 -- a ~1in strip on a blank page, invisible to make check (the key stranded the
% same box, so parity was perfect at 4/4). \workrowsep was swept 8pt->0pt and never
% bought the page back, so the 4th page is unavoidable; this \tcbbreak instead sends
% part (c) over with the remindbox so p4 carries substantial content. Mirrored in
% homework_key. Breaking one item earlier (before part (b)) was also tested: still
% 4 pages, but it left half of p3 empty, so this placement is the better balance.
% See references/conventions.md ("boxguard").
\tcbbreak
  \item Now change the $-2$ to $-\tfrac{1}{2}$ and nothing else. Name what changed about the shape,
        then find the two new inputs where $y=0$.
\begin{work}
  0 &= -\tfrac{1}{2}(x-4)^2+8 \\
  -8 &= -\tfrac{1}{2}(x-4)^2 \\
  16 &= (x-4)^2 \\
  x-4 &= \pm 4 \\
  x &= 0 \text{ or } x = 8
\end{work}

\writelines{2}
\end{enumerate}
\end{extensionbox}

\vspace{0.10in}

\begin{remindbox}
\small
\textbf{Next: Lesson 2.4 --- Intercepts, Zeros, Maximums and Minimums, and Where a Graph Rises and
Falls} (\texttt{AFDA.AF.2b--d}). Item 10 asked all three of those questions about the yearbook curve
in ordinary words. Next lesson gives each one its name and its notation --- and everything you name
will be read straight off a graph, which is exactly how this unit works.
\end{remindbox}

\end{document}
